\section{Conclusion and Future Work}
\label{submission:conclusion}

In this work, we deconstructed the state-of-the-art unsupervised test-time model merging framework, AdaMerging, through a highly critical, minimalist lens. Our theoretical and empirical investigations exposed two major optimization anomalies: the \textbf{Overfitting-Optimizer Paradox} and the \textbf{Spatial Averaging Paradox}.

First, we deconstructed the \textbf{Overfitting-Optimizer Paradox}: when given hundreds of layer-specific or parameter-specific degrees of freedom, test-time entropy minimization is prone to test-time overfitting. However, we demonstrated that the learned coefficients are structurally specialized and capture vital architectural hierarchy. Shuffling these coefficients across layers (\emph{Intra-Task Layer Shuffling}) collapses performance by disrupting this hierarchy, whereas post-hoc \emph{Spatial Averaging} smooths away the individual layer-wise overfitting component while successfully retaining the core, regularized task-level scaling signals. 

Second, we discovered and mathematically characterized the \textbf{Spatial Averaging Paradox}: while post-hoc spatial averaging succeeds, direct low-dimensional task-wise optimization fails spectacularly, actively degrading performance below its uniform initialization. We explained this behavior through multi-task gradient imbalance caused by uncalibrated prediction entropy objectives: under global task-wise constraints, the optimizer is heavily dominated by easy tasks, leading to destructive parameter interference and performance collapse on harder tasks.

Our findings advocate for a more rigorous, mathematically grounded approach to weight-space combinations. Occam's razor is a powerful guide, but we must understand the fundamental limits and uncalibrated objectives of test-time optimization to avoid pathologic optimization traps.

\subsection{Limitations and Future Directions}
While our deconstructive analysis provides deep insights into the optimization mechanics of adaptive model merging, we identify several boundary conditions and critical areas for future work:

\textbf{1. Test-Time Calibration Scale vs. Evaluation Rigor.}
While our test-time adaptation and calibration are performed on a small, sample-efficient subset of 64 unlabeled images per task to reflect realistic, data-scarce test-time conditions, our final evaluations are scaled to the full, standard test splits of each dataset (e.g., 10,000 images each for MNIST, FashionMNIST, and CIFAR-10, and 26,032 images for SVHN). This dual-scale protocol ensures that while the test-time adaptation remains highly practical and lightweight, our empirical findings are measured with maximum statistical rigor, yielding tight confidence intervals and eliminating evaluation split bias. Future work should investigate whether increasing the test-time calibration batch size itself could mitigate gradient imbalance.

\textbf{2. Boundary Conditions on Task Homogeneity.}
Our multi-task vision benchmark consists of highly heterogeneous tasks (ranging from greyscale digits to natural color images) with naturally unbalanced baseline prediction entropies, which triggered severe multi-task gradient imbalance. Investigating these paradoxes on homogeneous task setups (e.g., DomainNet or PACS, where task experts are fine-tuned on similar domains and exhibit naturally balanced entropy landscapes) would define the exact boundary conditions under which the Spatial Averaging Paradox manifests.

\textbf{3. Scaling to Large Language Models (LLMs).}
Finally, we aim to study how these weight-space optimization paradoxes scale to LLMs (e.g., merging instruction-tuned or domain-specialized language models). Under generative foundation models, test-time unsupervised adaptation is typically driven by minimizing token-level prediction perplexity or generation entropy. However, these objectives are highly sensitive to token-level gradient imbalance. In generative settings, highly frequent, easily predictable tokens (such as boilerplate templates, repetitive prompt phrases, or common function words) produce small loss values but massive, consistent gradients that dominate the joint perplexity loss landscape over rare, highly informative tokens associated with complex logical or reasoning tasks. This token-level and task-level imbalance directly mirrors the easy vs. hard task gradient imbalance observed in our vision models. To extend our paradoxes to generative language models, researchers should formulate balanced perplexity metrics or use calibration splits featuring curated, uniform in-context examples to avoid task-level gradient dominance. Translating our deconstructive insights and post-hoc spatial regularization to LLM weight-space combinations is a highly promising avenue for future research.

\textbf{4. Alternative Unsupervised Surrogate Losses.}
Since prediction entropy minimization is fundamentally a transductive, label-free objective prone to pathological logit-scaling shortcuts in low-dimensional bottlenecks, future research should explore more robust unsupervised surrogate losses. For example, self-supervised contrastive losses (e.g., CLIP InfoNCE) or mask-reconstruction objectives (e.g., Masked Autoencoder reconstruction loss) could be evaluated. Because these self-supervised objectives measure functional representational alignment rather than purely driving softmax sharpness, they are expected to bypass the uncalibrated nature of prediction entropy, potentially stabilizing direct flat task-wise optimization and mitigating weight-space interference.

\textbf{5. Generalization to Alternative Neural Network Families.}
Our empirical evaluations were conducted on the widely adopted CLIP ViT-B/32 Transformer backbone, where the hierarchical representational progression (early general visual features vs. late task-specific features) is well-documented. Future work should systematically evaluate whether the Overfitting-Optimizer and Spatial Averaging Paradoxes generalize to alternative neural network architectures. This includes hierarchical vision transformers like Swin Transformers \cite{swin} (where self-attention windows are shifted and resolutions vary across stages) and modern, fully convolutional networks like ConvNeXt \cite{convnext} (which emulate Transformer design choices using depthwise convolutions). Dissecting how structural routing behaves in non-isotropic and purely convolutional parameter spaces will provide a complete, architecture-agnostic validation of our deconstructive findings.

\textbf{6. Synergy with Static Weight-Space Sparsification.}
While this study contrasted adaptive model merging with static, non-adaptive baselines like TIES-Merging and DARE-Merging, these paradigms are fundamentally complementary rather than mutually exclusive. Static sparsification methods excel at removing fine-tuning noise and resolving sign conflicts, while adaptive scale regularization identifies optimal task contributions. A highly promising and exciting future research direction is to explore the direct synergy of these approaches: specifically, applying our post-hoc Spatial Averaging on top of pruned or sign-resolved base task vectors. Combining the noise-robustness of static coordinate-wise resolution with the unsupervised scaling properties of our spatial low-pass filter could yield even higher multi-task generalization and form a unified, state-of-the-art merging framework.
